%% Example CV - Marcus D. Yang
%% Comprehensive overview of all competencies, experience, and achievements.
%% Use as a master reference when tailoring role-specific CVs.
%%
%% Compile with: lualatex main_example.tex
%% (pdflatex often fails on modern MiKTeX with fontawesome5 font-expansion errors.)

\documentclass[11pt,a4paper,sans]{moderncv}
\moderncvstyle{banking}
\moderncvcolor{blue}

% Force the name and section headings to render in moderncv blue (color1).
% Default banking leaves them black: moderncvstylebanking.sty's \colorlet
% copies (not aliases) the pre-scheme accent colour, so the name colours are
% frozen before \moderncvcolor runs. Re-let them after. \namefont is the hook
% every name-style macro routes through, so this also works on moderncv 2.3.1
% (Debian/Ubuntu apt), which has no \firstnamestyle/\lastnamestyle at all.
\renewcommand*{\namefont}{\fontsize{34}{36}\bfseries\upshape}
\colorlet{firstnamecolor}{color1}
\colorlet{lastnamecolor}{color1}
\colorlet{namecolor}{color1}
\renewcommand*{\sectionstyle}[1]{{\sectionfont\color{color1}#1}}

\usepackage[utf8]{inputenc}
% moderncv loads hyperref itself in an \AtEndPreamble hook, so \hypersetup
% must go in an \AtEndPreamble of our own: on moderncv < 2.4 a top-level
% \usepackage{hyperref} clashes with the class's own
% \RequirePackage[unicode]{hyperref}. From 2.4.0 the class passes its options
% through \PassOptionsToPackage instead, which is what removes that clash.
\AtEndPreamble{\hypersetup{
    colorlinks=true,
    linkcolor=blue,
    filecolor=magenta,
    urlcolor=blue,
    pdftitle={Marcus D. Yang - CV},
    % Keep pdfpagemode=UseNone: this block runs after moderncv's own
    % \AtEndPreamble (moderncv.cls sets pdfpagemode there), so a FullScreen
    % value here would win and open every CV in fullscreen presentation mode.
    pdfpagemode=UseNone,
}}
\usepackage[scale=0.80]{geometry}
\usepackage{import}

% personal data
\name{Marcus}{Yang}
% If you have no address to list, DELETE this whole line. \address{}{}{} fails
% with "There's no line here to end" on every moderncv version.
\address{Champaign, IL, USA}{}{}
\phone[mobile]{+1 650 954 0288}
\email{yangdmarcus@gmail.com}
\extrainfo{\href{https://www.linkedin.com/in/marcusdyang}{LinkedIn}, \href{https://github.com/mdyang26}{GitHub}}

\begin{document}

\makecvtitle

% ============================================================
%     PROFILE STATEMENT
% ============================================================

\vspace{6pt}
\small{Aerospace Engineering student at the University of Illinois Urbana-Champaign (Computer Science minor, James Scholar Honors, GPA 3.57/4), in progress, expected May 2028. Hands-on experience across the full design-build-test cycle: CAD modeling and CFD analysis in Siemens NX and Ansys Fluent, embedded controls firmware in C++, and hardware fabrication and static-fire testing. Independently designed, built, and tested a two-axis thrust-vector-control rocket, root-causing two static-fire test failures directly from telemetry and video data. Seeking a paid, in-person internship in aerospace, mechanical, or manufacturing engineering.}

% ============================================================
%     CORE COMPETENCIES
% ============================================================

\section{Core Competencies}
\vspace{1pt}
\begin{itemize}

\item \textbf{CAD \& CFD}: Siemens NX (incl. NXOpen/Python scripting), Fusion 360, Ansys Fluent, XFLR5, OpenRocket -- full design cycle from geometry to aerodynamic validation

\item \textbf{Embedded Systems \& Controls}: C++ on Teensy/Arduino, closed-loop PID control, IMU sensor fusion, telemetry logging and analysis

\item \textbf{Manufacturing \& Fabrication}: CNC and laser cutting, composite layups (fiberglass/epoxy), FDM 3D printing and print-parameter tuning for structural parts, shop tools (spot welder, belt sander, band saw)

\item \textbf{Programming}: Python (NumPy, pandas, matplotlib, Monte Carlo simulation), C++, Java, JavaScript

\item \textbf{Test \& Data Analysis}: Root-cause failure analysis from flight/test telemetry, video, and audio data; signal analysis

\item \textbf{Electronics / PCB}: KiCad, Eagle PCB

\end{itemize}

% ============================================================
%     PROFESSIONAL EXPERIENCE
% ============================================================

\section{Professional Experience}
\vspace{3pt}
\begin{itemize}

% --- Independent Project ---
\needspace{5\baselineskip}
\item{\cventry{2026}{TVC Rocket -- Solo Design, Build \& Test}{Independent Project}{Champaign, IL}{}{\vspace{1pt}
\begin{itemize}
    \item {Designed, built, and tested a 4in-diameter model rocket with an active two-axis gimbaled thrust-vector-control system: CAD (Siemens NX), embedded firmware (Teensy 4.1, C++), closed-loop PID control, FDM 3D-printed structures}
    \item {Built a 6-DOF-informed flight simulation and ran a Monte Carlo sensitivity analysis (hundreds of runs) that found landing performance was far more sensitive to IMU accelerometer bias than servo speed, redirecting engineering effort toward sensor calibration}
    \item {Root-caused two independent static-fire test failures purely from telemetry: an axis-convention bug in the attitude estimator and an inverted control-sign bug, both diagnosed via statistical analysis of attitude logs and frame-by-frame video/audio tracking}
    \item {Logged telemetry at 83 Hz with zero dropouts across two static fires (15{,}728 and 12{,}611 samples)}
\end{itemize}}}

\vspace{3pt}

% --- Liquid Rocketry at Illinois ---
\needspace{5\baselineskip}
\item{\cventry{2025-Present}{Member}{Liquid Rocketry at Illinois}{Champaign, IL}{}{\vspace{1pt}
\begin{itemize}
    \item {Modeled P\&ID system architecture in Siemens NX and integrated components into the pressurant tank assembly}
    \item {Created CAD geometries for canard-grid fin configurations and ran Ansys Fluent CFD cases to quantify aerodynamic interactions across 0-20 deg deflection}
    \item {Engineered a KiCad PCB radio transmitter (USB-C + MCU) prototype for ground-test telemetry links}
\end{itemize}}}

\vspace{3pt}

% --- Flight Club Aerospace ---
\needspace{5\baselineskip}
\item{\cventry{2021-2024}{Landing Gear Research Lead}{Flight Club Aerospace}{}{}{\vspace{1pt}
\begin{itemize}
    \item {Manufactured wing and control surfaces (foam ribs, aluminum spars, epoxy/fiberglass layups) and aligned two 7-ft wings for flight testing}
    \item {Modeled gussets and ribs in Fusion 360; coordinated fit-checks between structural and aerodynamic components}
    \item {Presented landing-gear recommendations based on structural load and cost analyses}
\end{itemize}}}

\vspace{3pt}

% --- Asian American Association ---
\needspace{5\baselineskip}
\item{\cventry{2024-Present}{Treasurer | Fundraising Officer}{Asian American Association}{Champaign, IL}{}{\vspace{1pt}
\begin{itemize}
    \item {Managed membership operations for 200+ members, coordinating records and organizational logistics}
    \item {Secured partnerships with 15 on-campus sponsors to support events and fundraising initiatives}
\end{itemize}}}

\end{itemize}

% ============================================================
%     EDUCATION
% ============================================================

\section{Education}
\vspace{1pt}
\begin{itemize}

\needspace{5\baselineskip}
\item{\cventry{2024-2028}{B.S. in Aerospace Engineering, Computer Science Minor (James Scholar Honors)}{University of Illinois Urbana-Champaign}{Champaign, IL}{}{\vspace{1pt}
In progress, expected May 2028. GPA: 3.57/4. Relevant coursework: Incompressible Flow, Aerospace Control Systems, Aerospace Structures, Thermodynamics, Engineering Materials, Statics, Dynamics, Linear Algebra, Electrical and Electronic Circuits.
}}

\end{itemize}

% ============================================================
%     LANGUAGES
% ============================================================

\section{Languages}
\vspace{1pt}
\begin{itemize}
\item {English (native).}
\end{itemize}

% ============================================================
%     REFERENCES
% ============================================================

\section{References}
\vspace{1pt}
\begin{itemize}
\item{Available upon request.}
\end{itemize}

\end{document}
